% Imporditud: tools/import_eksperimendid.py
% Allikas: Eesti füüsikaolümpiaadi eksperimendi ülesannete
%   kogu (Rael Kalda, 2023), ülesanne Ü92, lahendus L92.
\setAuthor{EFO žürii}
\setRound{lõppvoor}
\setYear{2004}
\setNumber{P E1}
\setDifficulty{3}
\setTopic{dünaamika}
\setVahendid{5 juppi niiti, joonlaud, koormis massiga 100 g, statiiv vardaga, statiivi jalg}
\setPunktid{10}
\setAllikas{Eksperimendi kogu Ü92, lahendus lk 72}
\setLahendusseis{toores}

\prob{Niit}
Määrake suurim jõud, millega võib tõmmata antud niiti.

\hint

\solu
Üks võimalik lahendus põhineb kangi kasutamisel (vt. joon.). Joonlauda kasutada kangina, vardaga statiivi toetuspunktina, statiivi jalga niidi ühe otsa kinnitamiseks. Viis niiti on antud selleks, et teha viis mõõtmist ja leida keskmine tulemus.

$d_1$ $d_2$

koormis joonlaud

mg

T niit statiivi jalg
\probend
